\documentclass{article}
\usepackage{tcee-repaso}
\tcetema[El pensamiento económico de Keynes]{3.A.4}{El pensamiento económico de Keynes. Referencia a la economía postkeynesiana y el desequilibrio.}
\tcefecha{04/11/2026}
\begin{document}
\cabecerarepaso
\begin{columnas}

% ==========================================================================
\bloque{Esquema del cante}
\esquemacante{%
\apartadocante{Introducción: «¡tan largo me lo fiáis!»}{2,5}
\apartadocante{I. La escuela keynesiana y J.\,M.\,\autor{Keynes}}{6}
\subcante{1. Caracterización según \autor{Schumpeter}}{1,5}
\subcante{2. Contexto, método y obra}{2}
\subcante{3. El punto de partida: el modelo neoclásico}{2,5}
\apartadocante{II. Formalización y política económica}{10,5}
\subcante{1. Formalización (bienes 3,5; dinero 2; trabajo 2)}{7,5}
\subcante{2. Implicaciones de política económica}{2}
\subcante{3. Valoración: las lecturas de la \emph{Teoría general}}{1}
\apartadocante{III. Postkeynesianos y desequilibrio}{10}
\subcante{1. La escuela de Cambridge}{5,5}
\subcante{2. Los keynesianos del desequilibrio}{4,5}
\apartadocante{Conclusión}{1}}
\emph{Reparto}: \autor{Keynes} 16,5 min, herederos 10 (el título los deja como «referencia»). \emph{Pizarra} (intro y conclusión sin pizarra), 3 zonas. \emph{Izda.}: refrán y 6 pegatinas junto a su concepto: «¿Dónde va Vicente? Donde va la gente» (convenciones, cap.~12), «puedes llevar el caballo al río, pero no obligarlo a beber» (trampa de la liquidez), «desnudar a un santo para vestir a otro» (bajar salarios, cap.~19), «no se pueden sumar peras con manzanas» (controversia del capital), «después de la calma viene la tempestad» (\autor{Minsky}), «de donde no hay, no se puede sacar» (decisión dual). \emph{Centro}: gráficos, se borran al cambiar de bloque. \emph{Dcha.}: árbol de la escuela (se completa con los herederos) y, recuadradas hasta el final, demanda efectiva con multiplicador y preferencia por la liquidez. Anexo y profundización: para las preguntas.

% ==========================================================================
\bloque[2,5]{Introducción}
\begin{itemize}
\item \textbf{Enganche}: «¡Tan largo me lo fiáis!» (don Juan de Tirso de \autor{Molina}). Jul. 2020: Consejo Europeo, \emph{NextGenerationEU}, 750.000 M€ (precios 2018) con deuda común para sostener inversión y demanda; nadie propuso esperar a que bajaran los salarios. \autor{Keynes} (1923, \emph{Tract}): «a largo plazo, todos muertos». Pregunta: ¿se corrige sola una economía con paro, o hay que empujarla?
\item \textbf{Relevancia}: \emph{teoría}: nace la \concepto{macroeconomía} (producción, empleo y precios en conjunto; agregados de la contabilidad nacional \vertema{3.A.30}); \emph{política}: estabilización fiscal, Estado responsable del empleo; \emph{actualidad}: 2008-09 y 2020: multiplicadores fiscales (\autor{Blanchard}-\autor{Leigh}, 2013) \vertema{3.A.40}, tipos en cero \vertema{3.A.37}, \autor{Minsky} \vertema{3.A.42}; 2 de los 4 motivos de \autor{Schumpeter} (1954).
\item \textbf{Contextualización}: paradigmas y anomalías (\autor{Kuhn}, 1962) \vertema{3.A.1}. Clásicos y neoclásicos \vertema{3.A.2} \vertema{3.A.3}: ley de \autor{Say}, precios y salarios flexibles, dinero solo fija $P$. Anomalía: paro de entreguerras (R.\,Unido $\geq$9,7\,\% en 1921-29, 22,1\,\% en 1932; EE.\,UU. PIB real $-$26\,\% 1929-33, paro 25\,\%); bajar salarios no curaba. \autor{Keynes} (1936): \concepto{demanda efectiva} y \concepto{paro involuntario}. 3 lecturas: síntesis (IS-LM) \vertema{3.A.5}, postkeynesianos de Cambridge, desequilibrio $\to$ nueva economía keynesiana \vertema{3.A.7}.
\item \textbf{Problemática}: ¿qué es la demanda efectiva y por qué hay paro aunque bajen los salarios? ¿Papel del dinero, el tipo y la incertidumbre? ¿Qué política? ¿Era la síntesis la única lectura?
\item \textbf{Estructura}: \hyperref[rep:I]{I}, escuela, autor y modelo neoclásico de partida; \hyperref[rep:II]{II}, formalización (bienes, dinero, trabajo), política y lecturas; \hyperref[rep:III]{III}, Cambridge (\autor{Kalecki}, \autor{Robinson}, \autor{Kaldor}, \autor{Minsky}) y desequilibrio (\autor{Patinkin}, \autor{Clower}, \autor{Leijonhufvud}, \autor{Malinvaud}).
\end{itemize}

% ==========================================================================
\bloque[6]{I. La escuela keynesiana y J.\,M.\,\autor{Keynes}\label{rep:I}}
\begin{ideaclave}
\autor{Keynes} funda una escuela para explicar el paro involuntario como fallo del sistema en su conjunto, no de un mercado; lo hace con el paro británico de entreguerras delante y frente a un modelo neoclásico donde la flexibilidad asegura el pleno empleo y el dinero es neutral.
\end{ideaclave}

\sub{I.1. Caracterización según \autor{Schumpeter}\label{rep:I.1}}
\begin{itemize}
\item 3 rasgos: líder, seguidores, cuerpo doctrinal. \emph{Líder}: J.\,M.\,\autor{Keynes} (1883-1946).
\item \emph{Seguidores}, 3 ramas: \emph{síntesis} (\autor{Hicks}, \autor{Hansen}, \autor{Samuelson}, \autor{Modigliani}, \autor{Tobin}, \autor{Solow}, \autor{Klein}, \autor{Meade}) \vertema{3.A.5}; \emph{postkeynesianos} (\autor{Kahn}, J.\,\autor{Robinson}, \autor{Kaldor}, \autor{Sraffa}, \autor{Pasinetti}; precursor \autor{Kalecki}; EE.\,UU.: \autor{Minsky}, \autor{Davidson}, \autor{Weintraub}); \emph{desequilibrio} (antecedente \autor{Patinkin}; \autor{Clower}, \autor{Leijonhufvud}, \autor{Barro}, \autor{Grossman}, \autor{Malinvaud}).
\item \emph{Objeto}: probar que puede haber \concepto{paro involuntario} (quien quiere trabajar al salario vigente y no encuentra empleo) por demanda agregada insuficiente. Cap.~2: los clásicos explican el \emph{friccional} y el \emph{voluntario}, no el involuntario.
\item \emph{Método}: agregado y de CP; instrumental de \autor{Marshall} (competencia perfecta, razonamiento verbal), no \autor{Walras}; incertidumbre y expectativas centrales; ciencia moral al servicio de la política.
\item \emph{Supuestos frente a los neoclásicos}: horizonte CP («todos muertos»); \emph{incertidumbre} no probabilizable (cap.~12); salario \emph{nominal} negociado, se resiste a bajar (cap.~2) y bajarlo no garantiza el pleno empleo (cap.~19); precios también flexibles y en competencia perfecta; tipo de interés \emph{monetario} (cap.~13); dinero también depósito de valor; ley de \autor{Say} no se cumple (cap.~3).
\item Mantiene la \emph{competencia perfecta} (no la imperfecta de \autor{Robinson} y \autor{Chamberlin}, 1933 \vertema{3.A.18}): buscaba un fallo más profundo que el monopolio o los sindicatos.
\end{itemize}

\sub{I.2. Contexto, método y obra\label{rep:I.2}}
\begin{itemize}
\item \emph{Formación}: Cambridge 1883, hijo de J.\,N.\,\autor{Keynes}; matemáticas en King's (1905); economía con \autor{Marshall} y \autor{Pigou}. De \autor{Marshall}, el \emph{método}: CP, equilibrio en tiempo real, expectativas \vertema{3.A.16}.
\item 1906 India Office $\to$ \emph{Indian Currency and Finance} (1913); 1908 profesor en Cambridge; \emph{A Treatise on Probability} (1921): probabilidad como grado de creencia razonable, no siempre medible: raíz de la \emph{incertidumbre} \vertema{3.A.10}.
\item \emph{3 etapas} (\autor{Roncaglia}): (1) \emph{I Guerra Mundial}: Tesoro; París; \emph{The Economic Consequences of the Peace} (1919) contra las reparaciones; \emph{Tract} (1923): estabilidad de precios vs. oro a la paridad de preguerra; \emph{The End of Laissez-Faire} (1926). (2) \emph{Gran Depresión}: \emph{Treatise on Money} (1930), aún cuantitativo; «círculo» de Cambridge 1930-31 (\autor{Kahn}, J. y A.\,\autor{Robinson}, \autor{Sraffa}, \autor{Meade}) $\to$ multiplicador de \autor{Kahn} (1931) y \emph{Teoría general} (1936); réplica 1937. (3) \emph{II Guerra Mundial}: \emph{How to Pay for the War} (1940, ahorro obligatorio); Bretton Woods 1944: \emph{Plan Keynes} (unión de compensación, \emph{bancor}, ajuste simétrico) vs. \emph{Plan White}, que gana \vertema{3.B.20}; Savannah, mar. 1946 \vertema{3.B.22}; muere 21 abr. 1946.
\item \emph{Contexto}: 1921, PIB real R.\,Unido $-$9,7\,\%, paro asegurados 16,9\,\%; oro a la paridad de preguerra (abr. 1925) obliga a bajar salarios y precios; abandona el oro (sept. 1931); paro 22,1\,\% (1932).
\item \emph{Propósito}: no revolucionario: salvar el capitalismo de sí mismo; Estado responsable del empleo, con libertad y eficiencia de mercado (\autor{Skidelsky}).
\end{itemize}

\sub{I.3. El punto de partida: el modelo neoclásico\label{rep:I.3}}
\begin{itemize}
\item Reconstrucción (\autor{Marshall}, \autor{Pigou}, \autor{Fisher}, \autor{Wicksell}; \autor{Heijdra}, 2017). \emph{Supuestos}: economía cerrada, $K$ dado, competencia perfecta; precios, salarios y tipo \emph{flexibles}; sin ilusión monetaria; dinero solo medio de cambio. Causalidad: trabajo $\to$ fondos prestables $\to$ dinero.
\item \emph{Trabajo fija empleo} \vertema{3.A.25}:
\begin{formula}
\centering $w\equiv\dfrac{W}{P}=\mathrm{PMg}_L(L),\ \ L^D(w^c)=L^S(w^c)\Rightarrow L^c,\ \ Y_F=F(L^c)$
\end{formula}
\item[] OA \emph{vertical} en $Y_F$.
\item \emph{Bienes fija $r$}: \concepto{fondos prestables}, $S(r)=I(r)+G-T\Rightarrow r^n$ (tipo natural de \autor{Wicksell}; expresión de \autor{Robertson} y \autor{Ohlin}).
\item \emph{Dinero fija $P$}: \concepto{teoría cuantitativa} \vertema{3.A.35}: $MV=PY$ (\autor{Fisher}, 1911) o $M^D=kPY$ (Cambridge, \autor{Pigou}, 1917), $k=1/V$:
\begin{formula}
\centering $P=\dfrac{MV}{Y_F}=\dfrac{M}{kY_F}$
\end{formula}
\item \emph{Implicaciones}: ley de \autor{Say}; fiscal: \concepto{expulsión} total (cambia la composición, no el nivel) $\Rightarrow$ presupuesto equilibrado; dinero \concepto{neutral}, \concepto{dicotomía clásica}; solo paro friccional y voluntario; si persiste, sindicatos (\autor{Hicks}, 1932 \vertema{3.A.27}) o monopolio (\autor{Chamberlin}, 1933): bajar salarios, \emph{laissez-faire}.
\item \emph{Valoración}: coherente, pero no explicaba un paro masivo y duradero: la anomalía. De $r$ pasamos a $i$: el interés, precio de la liquidez y no del ahorro.
\end{itemize}

% ==========================================================================
\bloque[10,5]{II. Formalización y política económica\label{rep:II}}
\begin{ideaclave}
\autor{Keynes} invierte la causalidad neoclásica: la demanda efectiva fija producción y empleo; el tipo se fija en el mercado de dinero; el ahorro se ajusta a la inversión vía renta; hay paro involuntario aunque bajen los salarios. Remedio: política fiscal.
\end{ideaclave}

\sub{II.1. Formalización y comparación con el modelo neoclásico\label{rep:II.1}}
\begin{itemize}
\item Orden de la \emph{Teoría general}: bienes (demanda efectiva $\to$ $Y$) $\to$ dinero ($i$ $\to$ $I$) $\to$ trabajo (el empleo sale de los 2). Fórmulas = forma lineal posterior, no de \autor{Keynes}.
\item \textbf{Bienes} (cap.~3, el corazón). \emph{Oferta agregada} $Z(L)$: ingreso que compensa dar el empleo $L$; \emph{demanda agregada} $D(L)$: ingreso esperado. \concepto{Punto de demanda efectiva}: $D(L)=Z(L)$; nada asegura el pleno empleo. Ley de \autor{Say} $\equiv$ $D(L)=Z(L)$ $\forall L$.
\item \emph{Consumo} \vertema{3.A.33}: «función bastante estable» de la renta; \concepto{ley psicológica fundamental}: si la renta sube, el consumo sube, pero menos:
\begin{formula}
\centering $C_t=\overline C+c\,Y^d_t,\quad 0<c<1,\quad Y^d_t=(1-t)Y_t$
\end{formula}
\item[] PMeC $>$ PMgC y ↓ con la renta; ahorro $S=Y^d-C$ depende de la renta, no del tipo. $c$ constante = simplificación de la síntesis: para \autor{Keynes}, la PMgC tiende a ↓ con el empleo (cap.~10).
\item \emph{Inversión} \vertema{3.A.34}: \concepto{eficiencia marginal del capital} (EMgK), tipo de descuento que iguala rendimientos esperados y precio de oferta:
\begin{formula}
\centering $P^K=\displaystyle\sum_{t=1}^{T}\frac{R^e_t}{(1+\mathrm{EMgK})^t},\qquad I=\overline I-b\,i\ \ (b \text{ pequeño})$
\end{formula}
\item[] Se invierte hasta EMgK $=i$. \emph{Incertidumbre} (cap.~12): \emph{convenciones}; concurso de belleza; si la convención se rompe, inversión inestable: \concepto{\emph{animal spirits}}, «impulso espontáneo a la acción». \autor{Keynes} (1937): la incertidumbre, mensaje central \vertema{3.A.10}.
\item \emph{Equilibrio y} \concepto{multiplicador} (\autor{Kahn}, 1931, obras públicas; \autor{Keynes}, cap.~10, inversión): $Y=\mathrm{DA}$ a tipo dado:
\funcionclave{Y^*=\frac{\overline C+\overline I-b\,i_0+\overline G}{1-c(1-t)}\ \Rightarrow\ \Delta Y^*=\frac{\Delta\overline G}{1-c(1-t)}}
\item[] Multiplicador $>1$ (serie geométrica de rondas de gasto). Con $i$ variable, la misma ecuación da la IS \vertema{3.A.5}.
\end{itemize}
\graficorepaso[0.95\linewidth]{cruz-keynesiana}{Cruz keynesiana (\autor{Samuelson}, 1948): a la dcha. de $Y^*$, existencias ↑ (realizada $>$ planeada), $Y$↓}
\begin{itemize}
\item[] Se ajusta la \emph{renta}, no el tipo: el ahorro se iguala a la inversión.
\item \emph{Implicaciones}: no se cumple \autor{Say}; si $I$↓, $Y$ cae varias veces más. \concepto{Paradoja del ahorro}: si todos ahorran más, $Y$↓ y el ahorro nacional $S+(T-\overline G)$ sigue igual a $\overline I-bi_0$ (baja si $I$ depende de $Y$); \emph{Fábula de las abejas} (\autor{Mandeville}, cap.~23).
\item \textbf{Dinero}: el interés, recompensa por renunciar a la \emph{liquidez}, no al consumo (cap.~13). \emph{Supuestos}: dinero y bonos a LP; $M$ \emph{exógena}; cada inversor con su tipo «normal». \emph{3 motivos} \vertema{3.A.35}: transacción y precaución, $L_1(Y)$; \emph{especulación}, $L_2(i)$: si $i<$ normal, espera que suba (bonos bajan) y guarda dinero:
\funcionclave{\frac{M}{P}=L_1(Y)+L_2(i)=kY-h\,i,\quad k,h>0}
\item[] $\Rightarrow i=\tfrac1h(kY-M/P)$: fenómeno \emph{monetario}. $h=0$: Cambridge (clásico); $h\to\infty$: \concepto{trampa de la liquidez}, todo dinero de más se atesora (cap.~15; \autor{Keynes} no conocía ningún caso).
\item \emph{Implicaciones}: ataque a \autor{Say}: guardar dinero no es irracional; \autor{Say} (\autor{Lange}, 1942): valor de los excesos de demanda de bienes = 0; en economía monetaria, exceso de oferta de bienes con exceso de demanda de dinero (\autor{Becker}-\autor{Baumol}, 1952) \vertema{3.A.21}. Mercantilistas reivindicados (cap.~23). Dinero no neutral a CP; 2 eslabones débiles: $h$ grande y $b$ pequeño; demanda de dinero \emph{inestable}.
\item \textbf{Trabajo}: acepta el \emph{1.\textsuperscript{er} postulado} ($w=\mathrm{PMg}_L$), rechaza el \emph{2.º} (salario real = desutilidad marginal). Se negocia $W$, no $W/P$; resistencia a bajar el nominal (salario relativo), no a que suban precios: $W=W_0$.
\item[] Empleo lo fija la demanda; $\mathrm{PMg}_L=W_0/P$. Caída de DA: $Y$, $P$↓; demanda de trabajo $P\cdot\mathrm{PMg}_L$↓; con $W_0$, $L_0\to L_1$: ajuste por cantidades.
\item \concepto{Paro involuntario} (cap.~2) \vertema{3.A.28}: si ante una pequeña subida de precios respecto de $W$ la oferta y la demanda de trabajo superarían el empleo existente. OA creciente con recursos ociosos, vertical en $Y_F$.
\item \emph{¿Y si bajan los salarios?} (cap.~19, levanta la rigidez): solo crea empleo si sube $I$: tipo ↓ ($P$↓ $\Rightarrow$ $M/P$↑: \concepto{efecto \autor{Keynes}}) o EMgK ↑. Nada lo garantiza: trampa, deflación empeora expectativas y carga real de deudas; quita renta a quienes más consumen: «desnudar a un santo para vestir a otro». Salarios nominales estables $\Rightarrow$ más estabilidad.
\item \emph{Recapitulación}: en bienes, la renta iguala $S$ e $I$; en dinero, el tipo es monetario; en trabajo, el empleo lo fija la demanda.
\end{itemize}

\sub{II.2. Implicaciones de política económica\label{rep:II.2}}
\begin{itemize}
\item \emph{Fiscal}, la más segura: entra en la demanda y se multiplica; con paro no expulsa; déficit en depresión = remedio; obra pública aunque su utilidad directa sea escasa (cap.~10).
\item \emph{Monetaria}: útil pero insegura (2 eslabones débiles) \vertema{3.A.34}.
\item \emph{Comparación}: más gasto: neoclásicos ineficaz (expulsión) / \autor{Keynes} eficaz; más dinero: neutral / insegura; bajar salarios: eficaz / ineficaz (cap.~19); recomendación: \emph{laissez-faire} y presupuesto equilibrado / gestión de la demanda.
\item \emph{Más allá de la coyuntura} (cap.~24): tipo bajo permanente, «eutanasia del rentista»; «socialización bastante completa de la inversión» (no nacionalizar: el Estado asegura su volumen); impuestos progresivos (↑ propensión a consumir).
\item OA: demanda mueve $Y$ con recursos ociosos; en $Y_F$, solo precios \vertema{3.A.32}. Falta una teoría de la \emph{inflación} $\to$ \autor{Phillips} \vertema{3.A.5} \vertema{3.A.28}.
\end{itemize}

\sub{II.3. Valoración: las lecturas de la \emph{Teoría general}\label{rep:II.3}}
\begin{itemize}
\item \emph{Virtudes}: explica el paro masivo como fallo del sistema; funda la macroeconomía; rompe la dicotomía a CP. \autor{Samuelson} (1946): contagió a los menores de 35 «como una enfermedad\ldots{} a una tribu aislada de los mares del Sur».
\item \emph{Críticas}: demostración incompleta (quería paro sin rigideces y se leyó como salarios rígidos; \autor{De Vroey}, 2016); estático y de CP (sin inflación ni crecimiento $\to$ \autor{Harrod}); sin microfundamentos $\to$ desequilibrio y \vertema{3.A.7}.
\item \emph{\autor{Hicks}} (1937): IS-LM, $Y$ e $i$ a la vez \vertema{3.A.5}; \emph{\autor{Modigliani}} (1944): con salarios flexibles y sin trampa, pleno empleo $\Rightarrow$ \autor{Keynes} = caso particular de salarios rígidos.
\item \emph{¿Traducción o traición?} \autor{Patinkin} (1987): voz de \autor{Marshall}, razonamiento de \autor{Walras}; la síntesis se quedó con lo 2.º.
\end{itemize}

% ==========================================================================
\bloque[10]{III. Postkeynesianos y desequilibrio\label{rep:III}}
\begin{ideaclave}
Postkeynesianos: rechazan el instrumental neoclásico y llevan demanda efectiva e incertidumbre al LP, la distribución y las finanzas. Desequilibrio: aceptan ese instrumental y buscan en él los fundamentos de \autor{Keynes} (precios que no vacían, ajuste por cantidades). Los 2: el paro no se explica solo con un salario que no baja.
\end{ideaclave}

\sub{III.1. La economía postkeynesiana: la escuela de Cambridge\label{rep:III.1}}
\begin{itemize}
\item \emph{Líder}: J.\,\autor{Robinson}, con \autor{Kaldor}; núcleo, el «círculo» (\autor{Kahn}, \autor{Robinson}, \autor{Sraffa}); luego \autor{Kalecki}, \autor{Pasinetti}; EE.\,UU.: \autor{Minsky}, \autor{Davidson}, \autor{Weintraub} (\autor{King}, 2002). IS-LM = «\concepto{keynesianismo bastardo}» (\autor{Robinson}, 1962): olvida el tiempo histórico y la demanda efectiva a LP.
\item \emph{Objeto}, 4 ideas (\autor{Lavoie}, 2014): demanda efectiva también a LP ($I$ determina $S$); incertidumbre fundamental; distribución según instituciones y poder; dinero endógeno y finanzas inestables. \emph{Método}: tiempo histórico, clases sociales, precios con margen, sin función de producción agregada.
\item \textbf{\autor{Kalecki}}: demanda efectiva por su cuenta, desde \autor{Marx}, en polaco en 1933. Beneficios (cerrada, sin Estado, trabajadores no ahorran): $B=I+C_{\mathrm B}$; «los capitalistas ganan lo que gastan y los trabajadores gastan lo que ganan» (\autor{Kaldor}, 1956).
\item[] \concepto{Grado de monopolio}: margen $\mu$ sobre coste directo; $j$ = materias primas/salarios:
\begin{formula}
\centering $\dfrac{WL}{PY}=\dfrac{1}{1+\mu(1+j)},\qquad \mathcal L=\dfrac{\mu}{1+\mu}$ \ \vertema{3.A.17}
\end{formula}
\item[] «Paradoja de los costes»: subir salarios reales puede ↑ demanda, empleo y beneficios; crecimiento por salarios o beneficios (\autor{Bhaduri}-\autor{Marglin}, 1990) \vertema{3.A.45}. \emph{Ciclo político} (1943): el pleno empleo duradero da poder a los trabajadores y los empresarios presionan por la austeridad \vertema{3.A.23}.
\item \textbf{\autor{Harrod}} (1939): tasa \emph{garantizada} $s/v$ vs. \emph{natural}; «filo de la navaja» \vertema{3.A.43}. \autor{Solow}: $v$ variable; \autor{Kaldor} (1956): la \emph{distribución}, con $s_{\mathrm B}>s_{\mathrm L}$:
\begin{formula}
\centering $\dfrac{B}{Y}=\dfrac{1}{s_{\mathrm B}-s_{\mathrm L}}\dfrac{I}{Y}-\dfrac{s_{\mathrm L}}{s_{\mathrm B}-s_{\mathrm L}}$
\end{formula}
\item[] La inversión fija la parte de los beneficios. \autor{Pasinetti} (1962): aunque los trabajadores ahorren y tengan capital, \emph{ecuación de Cambridge} $B/K=\gamma_Y/s_{\mathrm B}$: no depende de tecnología ni del ahorro de los trabajadores \vertema{3.A.45}.
\item \textbf{Controversia del capital}: ¿en qué unidades se suman máquinas distintas? Solo en valor, que depende del tipo de beneficio (\autor{Robinson}, 1953): «no se pueden sumar peras con manzanas». \autor{Sraffa} (1960): fronteras salario-beneficio; si se cortan 2 veces, \concepto{retorno de técnicas} (\emph{reswitching}) y \emph{reversión del capital}. \autor{Levhari} (1965) erró; \autor{Samuelson} (1966) lo admitió. Distribución sin productividades marginales; la macro siguió con la función agregada \vertema{3.A.11} \vertema{3.A.43}.
\item \textbf{\autor{Minsky}}: inestabilidad nace dentro del sistema financiero (1975, 1986). \emph{3 posiciones} (1992): \emph{cobertura} (pagan intereses y principal), \emph{especulativa} (solo intereses; refinanciar principal), \emph{Ponzi} (ni intereses: solo si suben los precios de los activos). Expansión larga: cobertura $\to$ especulativa $\to$ Ponzi; «la estabilidad es desestabilizadora»; «después de la calma viene la tempestad». Remedio: «gran gobierno» y «gran banco» (prestamista de última instancia), pero animan el riesgo: regulación por delante \vertema{3.A.42}.
\item \emph{Dinero endógeno}: lo crean los bancos al prestar; el banco central fija el tipo, oferta «horizontal» (\autor{Kaldor}, 1982; \autor{Moore}, 1988) \vertema{3.A.37}.
\item \emph{Evidencia}: «momento \autor{Minsky}» (\autor{Cassidy}, 2008): deuda privada no financiera de EE.\,UU. 125\,\% del PIB (1990) $\to$ 171\,\% (fin 2007); hogares 61 $\to$ 98\,\%; máx. 172,0\,\% (3T 2008); 140,3\,\% (fin 2025). Banco de Inglaterra (2014): el dinero, depósitos que crean los bancos.
\item \emph{Valoración}: tomaron en serio incertidumbre, distribución y finanzas; anticiparon el dinero endógeno. Pero sin modelo común ni contraste sistemático, fuera de la corriente principal; la controversia, ganada en la lógica, no cambió la práctica.
\end{itemize}

\sub{III.2. Los keynesianos del desequilibrio\label{rep:III.2}}
\begin{itemize}
\item \emph{Líderes}: \autor{Clower} y \autor{Leijonhufvud} (finales de los 60): la \emph{Teoría general} trata de precios que \emph{no} vacían; el IS-LM la perdió. Seguidores: \autor{Patinkin} (antecedente), \autor{Barro}-\autor{Grossman}, \autor{Bénassy}, \autor{Drèze}, \autor{Malinvaud}; escuela de los 70, europea.
\item \emph{Objeto}: \concepto{equilibrios no walrasianos} (con racionamiento): precios lentos; se intercambia lo que quiere el \emph{lado corto}; el racionamiento se contagia: fallo de \emph{coordinación}. \emph{Método}: agentes que maximizan, sin subastador, con restricciones de cantidad.
\item \textbf{\autor{Patinkin}} (1956): dinero en el equilibrio general de \autor{Walras} \vertema{3.A.21}; \concepto{efecto saldos reales}: $Y^D=C(Y,i,M/P)+I(i)+\overline G$, $\partial C/\partial(M/P)>0$. Cae la demanda, precios lentos: empleo \emph{fuera} de las curvas de oferta y demanda de trabajo; paro involuntario sin salario real alto. No es equilibrio: deflación + saldos reales $\to$ pleno empleo, despacio; la política \emph{acorta} el ajuste. Críticas: paro transitorio, efecto pequeño y lento; \autor{Lucas} (1980): la versión más acabada de la síntesis.
\item \textbf{\autor{Clower}} (1965), «The Keynesian Counter-Revolution»: consumo según la renta \emph{realizada}. \concepto{Demanda nocional} (vende todo el trabajo deseado) vs. \concepto{demanda efectiva} $C^{\mathrm{ef}}=C(w\overline L)<C^{\mathrm{noc}}$: \concepto{decisión dual} («de donde no hay, no se puede sacar»). Ley de \autor{Walras} nocional:
\begin{formula}
\centering $\sum_k p_k z_k^{\mathrm{noc}}(\vec p)\equiv0;\quad \sum_{k\neq L}p_k z_k^{\mathrm{ef}}+W(\overline L-L^S)<0$
\end{formula}
\item[] El paro no envía señal al mercado de bienes: equilibrio con paro; el multiplicador es el contagio. Dilema: o \autor{Walras} es incompatible con \autor{Keynes} o \autor{Keynes} no aportaba nada. \autor{Clower} (1967): el dinero, antes que la compra. Críticas: en términos efectivos vuelve \autor{Walras}; no explica por qué los precios no se ajustan.
\end{itemize}
\graficorepaso[0.9\linewidth]{clower-decision-dual}{Decisión dual: con empleo racionado $\overline L$, la demanda efectiva cae bajo la nocional}
\begin{itemize}
\item \textbf{\autor{Leijonhufvud}} (1968): «economía keynesiana» (síntesis) vs. «economía de \autor{Keynes}». 4 tesis: cantidades se ajustan antes que precios (orden inverso a \autor{Marshall}); sin subastador, nadie conoce los precios de equilibrio; multiplicador vía decisión dual; tipo a LP que no baja (tipo «normal»): $S$ e $I$ no se coordinan en el tiempo. Paro = fallo de coordinación por información. \emph{Pasillo} (1973): pequeñas perturbaciones, neoclásicos; grandes (1929), \autor{Keynes}.
\item \textbf{\autor{Barro}-\autor{Grossman}} (1971): $P$ y $W$ fijos; empresas que no venden: $L^{\mathrm{ef}}=F^{-1}(Y)$; hogares: $C^{\mathrm{ef}}=C(\tfrac WP L^{\mathrm{ef}},\tfrac MP)$; $Y=C^{\mathrm{ef}}+\overline I+\overline G$: multiplicador keynesiano. Con exceso de demanda en los 2, multiplicador de oferta. \autor{Barro} (1979) lo abandona: precios fijos sin explicar $\to$ nueva macro clásica \vertema{3.A.6}.
\item \textbf{\autor{Malinvaud}} (1977): sindicatos vs. patronal en Francia: los 2 pueden tener razón.
\funcionclave{Y=\min\Big\{Y^D\big(\overline A,\tfrac MP\big),\ Y^S\big(\tfrac WP\big),\ F(\overline L)\Big\}}
\item[] \concepto{Paro keynesiano} ($Y^D$): exceso de oferta de bienes y trabajo $\Rightarrow$ más demanda (subir salarios puede ayudar: sindicatos). \concepto{Paro clásico} ($Y^S$): exceso de demanda de bienes y oferta de trabajo $\Rightarrow$ bajar salario real (patronal). \concepto{Inflación reprimida} ($F(\overline L)$): exceso de demanda en los 2 $\Rightarrow$ menos demanda. Se tocan en $\mathcal W$, único punto sin racionamiento. \emph{Subconsumo} (\autor{Muellbauer}-\autor{Portes}, 1978): fuera del modelo básico.
\end{itemize}
\graficorepaso[\linewidth]{malinvaud-regimenes}{Regímenes de \autor{Malinvaud} en $(W,P)$: K, C, I; frontera K-I horizontal}
\begin{itemize}
\item[] Evidencia: \autor{Drèze}-\autor{Bean} (1990), 10 países: paro europeo de los 80, parte clásica (costes tras el petróleo) y parte keynesiana \vertema{3.A.28}. Crítica: $W$ y $P$ exógenos, estático.
\item \emph{Valoración}: microfundamentos; paro como coordinación (equilibrios múltiples, \autor{Cooper}-\autor{John}, 1988) \vertema{3.A.21}; tipos de paro. Críticas: no explican la rigidez (\autor{Barro}, 1979); crítica de \autor{Lucas} \vertema{3.A.6}. \emph{Legado}: la nueva economía keynesiana explica la rigidez (costes de menú, salarios de eficiencia, contratos) \vertema{3.A.7}.
\item \emph{Comparación} (¿por qué paro?): \autor{Keynes}, demanda efectiva insuficiente; síntesis, salario nominal rígido; postkeynesianos, demanda efectiva también a LP; desequilibrio, precios que no vacían. Incertidumbre central para \autor{Keynes} y postkeynesianos; política: fiscal / sintonía fina / fiscal, rentas y regulación financiera / según el régimen.
\end{itemize}

% ==========================================================================
\bloque[1]{Conclusión}
\begin{itemize}
\item \textbf{Recapitulación}: «En conclusión,» \autor{Keynes} invierte la causalidad: demanda efectiva fija $Y$ y $L$; el interés, precio de la liquidez; bajar salarios no garantiza el pleno empleo. Cambridge lo lleva a distribución, capital y finanzas; el desequilibrio, a un fallo de coordinación.
\item \textbf{Relevancia}: la Gran Depresión le dio la razón; la estanflación, a sus críticos; Gran Recesión y pandemia: multiplicadores, trampa de la liquidez, «momento \autor{Minsky}».
\item \textbf{Extensiones}: síntesis y monetarismo \vertema{3.A.5}; nueva macro clásica \vertema{3.A.6}; nueva economía keynesiana \vertema{3.A.7}; dinero \vertema{3.A.35}; ciclo financiero \vertema{3.A.42}; crecimiento y distribución \vertema{3.A.43} \vertema{3.A.45}.
\item \textbf{Opinión}: el acierto fue la pregunta (¿paro como fallo del sistema?); hoy se responde con modelos que él no escribiría, pero la respuesta práctica de 2020 fue la suya.
\item \textbf{Idea final}: ¿se corrige sola? Quizá, pero esperar al LP es contestar «¡tan largo me lo fiáis!»; sus herederos añaden que ese LP puede no llegar nunca.
\end{itemize}

% ==========================================================================
\bloque{Autores y fechas}
\hito{1890}{\autor{Marshall}: \emph{Principles} (mercado de pescado)}
\hito{1911}{\autor{Fisher}: ecuación de cambio}
\hito{1917}{\autor{Pigou}: saldos de caja (Cambridge)}
\hito{1919}{\autor{Keynes}: \emph{Economic Consequences of the Peace}}
\hito{1921}{\autor{Keynes}: \emph{Treatise on Probability}}
\hito{1923}{\autor{Keynes}: \emph{Tract} («todos muertos»)}
\hito{1926}{\autor{Keynes}: \emph{The End of Laissez-Faire}}
\hito{1930}{\autor{Keynes}: \emph{Treatise on Money}}
\hito{1931}{\autor{Kahn}: multiplicador}
\hito{1933}{\autor{Kalecki}: demanda efectiva; \autor{Robinson}, \autor{Chamberlin}: competencia imperfecta}
\hito{1936}{\autor{Keynes}: \emph{Teoría general}}
\hito{1937}{\autor{Hicks}: IS-LM; \autor{Keynes}: réplica (incertidumbre)}
\hito{1939}{\autor{Harrod}: filo de la navaja}
\hito{1940}{\autor{Keynes}: \emph{How to Pay for the War}}
\hito{1943}{\autor{Kalecki}: ciclo político}
\hito{1944}{Bretton Woods; \autor{Modigliani}: salario rígido}
\hito{1946}{\autor{Keynes} muere (21 abr.); \autor{Samuelson}: la «enfermedad»}
\hito{1948}{\autor{Samuelson}: cruz keynesiana}
\hito{1953}{\autor{Robinson}: crítica del capital}
\hito{1956}{\autor{Patinkin}: \emph{Money, Interest, and Prices}; \autor{Kaldor}: distribución}
\hito{1960}{\autor{Sraffa}: \emph{Production of Commodities}}
\hito{1962}{\autor{Pasinetti}: ecuación de Cambridge; \autor{Robinson}: «keynesianismo bastardo»}
\hito{1965}{\autor{Clower}: decisión dual}
\hito{1966}{\autor{Samuelson}: admite el retorno de técnicas}
\hito{1968}{\autor{Leijonhufvud}: economía de \autor{Keynes}}
\hito{1971}{\autor{Barro}-\autor{Grossman}: desequilibrio general}
\hito{1977}{\autor{Malinvaud}: regímenes del paro}
\hito{1979}{\autor{Barro}: «Second Thoughts»}
\hito{1986}{\autor{Minsky}: \emph{Stabilizing an Unstable Economy}}
\hito{1990}{\autor{Drèze}-\autor{Bean}: paro europeo}

% ==========================================================================
\bloque{Datos clave}
\begin{itemize}
\item \emph{NextGenerationEU}: \dato{750.000 M€ (precios 2018)}{CE, 2020}.
\item R.\,Unido: PIB real 1921 \dato{$-$9,7\,\%}{BoE}; paro asegurados \dato{16,9\,\% (1921)}{BoE}, \dato{$\geq$9,7\,\% (1921-29)}{BoE}, \dato{22,1\,\% (1932)}{BoE}; sobre población activa \dato{15,4\,\% (1932)}{BoE}.
\item EE.\,UU.: PIB real \dato{$-$26\,\% (1929-33)}{BEA}; paro \dato{25\,\%}{Census, HSUS}.
\item Deuda privada no financiera EE.\,UU.: \dato{125\,\% $\to$ 171\,\% del PIB (1990-2007)}{BIS}; hogares \dato{61 $\to$ 98\,\%}{BIS}; máx. \dato{172,0\,\% (3T 2008)}{BIS}; \dato{140,3\,\% (4T 2025)}{BIS}.
\end{itemize}

% ==========================================================================
\bloque{Preguntas del tribunal}
\begin{enumerate}
\item \textbf{¿Por qué «general»?} Se ocupa del sistema en su conjunto (renta, empleo, ahorro agregados), no de una industria o un individuo (prefacio a la ed. francesa, 1939).
\item \textbf{¿Necesita \autor{Keynes} salarios rígidos?} No: cap.~19, bajarlos no cura el paro (trampa, expectativas, deudas, consumo). La rigidez es de la síntesis (\autor{Modigliani}, 1944).
\item \textbf{¿Qué fija el tipo de interés?} La preferencia por la liquidez y $M$, no ahorro e inversión; el ahorro se ajusta vía renta.
\item \textbf{¿Paradoja del ahorro?} Más ahorro deseado $\Rightarrow$ $Y$↓; el ahorro total no sube (o baja si $I$ depende de $Y$).
\item \textbf{¿Demanda nocional vs. efectiva?} Nocional: sin restricciones de cantidad; efectiva: con la renta realizada (\autor{Clower}).
\item \textbf{¿Cuándo tiene razón la patronal?} En paro clásico de \autor{Malinvaud}: salario real demasiado alto; en keynesiano, los sindicatos.
\item \textbf{¿Qué demostró el retorno de técnicas?} No se puede ordenar técnicas por intensidad de capital: la distribución no se explica por la productividad marginal.
\item \textbf{¿Por qué fracasó el desequilibrio?} No explicaba la rigidez de precios (\autor{Barro}, 1979) ni respondía a \autor{Lucas}; lo recoge la nueva economía keynesiana.
\end{enumerate}

\begin{trampa}
\begin{itemize}
\item \autor{Keynes} no supuso la PMgC constante (decreciente con el empleo): fue la síntesis.
\item \autor{Keynes} mantiene la \emph{competencia perfecta} y los precios flexibles; lo rígido es el salario \emph{nominal}.
\item IS-LM es de \autor{Hicks} (1937), no de \autor{Keynes}; cruz keynesiana: \autor{Samuelson} (1948); multiplicador: \autor{Kahn} (1931).
\item Trampa de la liquidez: demanda de dinero \emph{horizontal} ($h\to\infty$); \autor{Keynes} no conocía ningún caso. $h=0$: Cambridge.
\item Cruz: a la dcha. de $Y^*$, inversión realizada \emph{mayor} que la planeada (acumulación).
\item Para \autor{Keynes}, «clásicos» incluye a \autor{Marshall} y \autor{Pigou}; «neoclásico», término de \autor{Veblen} (h. 1900); «fondos prestables», de \autor{Robertson} y \autor{Ohlin}.
\item «Demanda efectiva», 3 sentidos: \autor{Keynes} (punto, cap.~3), \autor{Clower} (restricción de cantidad), \autor{Chamberlin} (curva de cuota) \vertema{3.A.18}.
\item Paro clásico de \autor{Malinvaud}: exceso de \emph{demanda} de bienes; subconsumo no es régimen del modelo básico.
\item Bretton Woods: ganó el \emph{Plan White}. Teoría cuantitativa: \autor{Azpilcueta} (1556) antes que \autor{Bodin} (1568).
\end{itemize}
\end{trampa}

\end{columnas}
\end{document}
